\chapter{Future Work}

This chapter outlines follow-up work in priority order, with the expected impact
of each item tied back to the limitations already identified.

The priorities are linked directly to the limitations in
Table~\ref{tab:limitations-impact-mitigation}. Items that strengthen every
quantitative claim are ranked above items that broaden scope or improve
deployment readiness.

\section{Priority Roadmap}

\begin{enumerate}
\item \textbf{Extend commit-locked multi-trial coverage to all supporting targets.}
The latest headline runs already use a clean-tree \(N=5\) policy, but older
supporting targets such as the unoptimised LowMC baseline and AES backend variants
should be rerun under the same policy. The output should include variability
bands, outlier notes, and regenerated figures from the same locked commit.

\item \textbf{Add per-target commit hashes to harness artifacts.}
Embedding git revision metadata directly in trial and aggregate JSON would remove
forensic ambiguity when rerunning months later. This should include dirty-tree
status so a successful run cannot be mistaken for a clean release build.

\item \textbf{Execute explicit AES backend feature-flag benchmarks.}
A strict \texttt{sbox-table} vs \texttt{sbox-logic} study would resolve whether
the optimisation workspace can produce measurable zkVM gains.

\item \textbf{Expand cross-zkVM replication.}
Porting the same workload contract to at least one additional zkVM would test
which conclusions are stack-specific and which generalize. The port should keep
input sizes, transcript framing, and benchmark IDs as close as possible to the
RISC Zero campaign.

\item \textbf{Extend transport-layer authentication.}
Adding signature-based sender binding for proof artifacts would close the current
deployment gap between data authenticity and sender identity. Replay controls and
session identifiers should be included so old valid receipts cannot be presented
as fresh protocol messages.

\item \textbf{Add negative-path authenticated-mode tests.}
Explicit tests for AAD, IV, ciphertext, tag, key, and image-ID mismatches would
turn the design-level failure analysis into executable regression coverage.

\item \textbf{Tighten the security argument into a fuller reduction chain.}
Turning the Chapter 10 game transitions into theorem-level reductions would
strengthen the cryptographic assurance story. This should separate confidentiality
terms from integrity, soundness, and replay assumptions.
\end{enumerate}

\section{Expected Impact}

Items 1 and 2 have the highest impact on report quality because they improve the
credibility and auditability of the remaining mixed-provenance quantitative
claims.

Items 3 and 4 determine how far findings can be generalized beyond the current
stack.

Items 5, 6, and 7 strengthen deployment readiness, negative-path assurance, and
formal-assurance depth, making the pipeline more suitable for real-world
integration and external review.

\section{Limitations to Future Work Mapping}

Table~\ref{tab:future-work-limitations-map} maps each stated limitation to the
future-work item that directly mitigates it.

\begin{table}[H]
\centering
\reporttablesetup
\begin{tabularx}{\textwidth}{@{}LLL@{}}
\hline
Limitation addressed & Future work item & Priority reason \\
\hline
Mixed-provenance support snapshots & Commit-locked multi-trial rerun for all
supporting targets & Highest impact on confidence in remaining legacy
comparisons \\
Incomplete artifact provenance & Per-target commit metadata & Required for
forensic reproducibility and external audit \\
Unresolved AES backend question & Feature-flag backend benchmark & Clarifies
whether AES optimisation work has measurable zkVM value \\
Single-zkVM scope & Cross-zkVM replication & Tests whether the AES/LowMC ranking
is stack-specific \\
Transport identity gap & Signature and replay-control layer & Moves the system
closer to deployable protocol design \\
Negative-path assurance gap & Authenticated-mode mismatch tests & Checks that
failure cases remain enforced as code evolves \\
Proof-sketch depth & Fuller reduction chain & Raises assurance from structured
argument to more formal cryptographic reasoning \\
\hline
\end{tabularx}
\caption{Mapping from stated limitations to prioritized future work.}
\label{tab:future-work-limitations-map}
\end{table}
